\documentclass[11pt]{article}

\usepackage[utf8]{inputenc}
\usepackage[T1]{fontenc}
\usepackage{amsmath,amssymb}
\usepackage{booktabs}
\usepackage{geometry}
\usepackage{graphicx}
\usepackage{float}
\usepackage[hidelinks]{hyperref}

\geometry{margin=1in}

\title{Modelo simplificado de reputación con tipo oculto}
\author{}
\date{}

\begin{document}

\maketitle

\section{Objetivo}

Este documento resume el modelo simplificado de reputación con producto oculto
implementado en \texttt{simulate\_hidden\_type\_model.py}. El objetivo es
reemplazar el estado Beta posterior del modelo original por un estado escalar
más transparente:
\[
\mu_t
= \mathbb{P}\{\text{Seller A es tipo alto / producto 2}\mid \mathcal{H}_t\}.
\]
El usuario conoce las probabilidades de éxito de los dos productos, pero no
observa qué producto usa el vendedor. Por lo tanto, la reputación relevante del
vendedor es la creencia \(\mu_t\), y el objeto central es la función de valor
\(V_{p_0}(\mu)\), que depende de la opción externa \(p_0\).

\section{Primitivas}

Seller A puede elegir un producto \(x\in\{1,2\}\). El producto 1 tiene menor
probabilidad de éxito y menor costo, mientras que el producto 2 tiene mayor
probabilidad de éxito y mayor costo:
\[
0<p_1<p_2<1,
\qquad
c_1<c_2.
\]
Si el usuario elige a Seller A, el vendedor recibe ingreso \(R\), paga el costo
\(c_x\), y se observa un resultado binario \(Y_t\in\{0,1\}\) con
\[
\mathbb{P}(Y_t=1\mid x)=p_x.
\]
La opción externa tiene probabilidad conocida de éxito \(p_0\). Los parámetros
base usados en los numéricos son:

\begin{table}[H]
\centering
\begin{tabular}{lc}
\toprule
Parámetro & Valor \\
\midrule
\(p_1\) & \(0.35\) \\
\(p_2\) & \(0.80\) \\
\(c_1\) & \(0.05\) \\
\(c_2\) & \(0.65\) \\
\(R\) & \(1.00\) \\
\(\gamma\) & \(0.98\) \\
\(\mu_{\mathrm{init}}\) & \(0.50\) \\
\texttt{grid\_size} & \(1000\) \\
\texttt{epsilon} & \(10^{-5}\) \\
\texttt{tol} & \(10^{-9}\) \\
\bottomrule
\end{tabular}
\caption{Parámetros base.}
\label{tab:parameters}
\end{table}

La probabilidad de éxito percibida de Seller A es
\[
\bar p(\mu)=p_1+\mu(p_2-p_1).
\]
El umbral de indiferencia del usuario en el modo de demanda determinista es
\[
\mu_0(p_0)=\frac{p_0-p_1}{p_2-p_1}.
\]
Con los parámetros base, \(p_2-p_1=0.45\). Por ejemplo,
\(\mu_0(0.50)=1/3\) y \(\mu_0(0.65)=2/3\).

\section{Demanda del usuario}

Se implementan dos modos de demanda.

\subsection{Demanda por umbral de media}

En el modo \texttt{mean\_threshold}, el usuario compara la calidad percibida
\(\bar p(\mu)\) con la opción externa \(p_0\). La demanda es
\[
D(\mu;p_0)=
\begin{cases}
1, & p_0\le p_1,\\
\mathbf{1}\{\mu\ge \mu_0(p_0)\}, & p_1<p_0\le p_2,\\
0, & p_0>p_2.
\end{cases}
\]
Este modo genera un corte duro. Si la creencia inicial queda por debajo del
umbral, no hay demanda, no se observan resultados y la creencia permanece fija.
Por eso, por ejemplo, con \(p_0=0.65\) y \(\mu_{\mathrm{init}}=0.5\), el usuario
no elige al vendedor: \(\bar p(0.5)=0.575<0.65\).

\subsection{Thompson sampling sobre el tipo oculto}

En el modo \texttt{ts\_type}, el usuario samplea un tipo
\(\tilde \omega\sim\mathrm{Bernoulli}(\mu)\). Si samplea el tipo alto, compara
\(p_2\) con \(p_0\); si samplea el tipo bajo, compara \(p_1\) con \(p_0\). La
demanda esperada es
\[
D(\mu;p_0)=
\begin{cases}
1, & p_0\le p_1,\\
\mu, & p_1<p_0\le p_2,\\
0, & p_0>p_2.
\end{cases}
\]
Dentro de la región intermedia \(p_1<p_0\le p_2\), \(p_0\) ya no entra en la
ecuación de Bellman: la demanda esperada es simplemente \(\mu\). Por eso las
políticas de \texttt{ts\_type} son casi invariantes en \(p_0\) dentro de esa
región; las diferencias en simulación vienen principalmente de ruido Monte
Carlo.

\section{Actualización bayesiana}

Si el usuario elige a Seller A y observa un éxito, la creencia se actualiza a
\[
\mu^S
=
\frac{\mu p_2}
{\mu p_2+(1-\mu)p_1}.
\]
Si observa un fracaso, la creencia se actualiza a
\[
\mu^F
=
\frac{\mu(1-p_2)}
{\mu(1-p_2)+(1-\mu)(1-p_1)}.
\]
Si el usuario no elige a Seller A, no hay nueva información y la creencia queda
igual.

\section{Problema del vendedor}

El vendedor observa \(\mu\) y elige \(x\in\{1,2\}\). La ecuación de Bellman para
un \(p_0\) fijo es
\[
\begin{aligned}
V_{p_0}(\mu)
&=
\gamma[1-D(\mu;p_0)]V_{p_0}(\mu) \\
&\quad
+D(\mu;p_0)\max_{x\in\{1,2\}}
\left\{
R-c_x
+\gamma\left[
p_xV_{p_0}(\mu^S)
+(1-p_x)V_{p_0}(\mu^F)
\right]
\right\}.
\end{aligned}
\]
El producto 2 es óptimo si el valor reputacional de mover la creencia hacia
\(\mu^S\) en vez de \(\mu^F\) compensa el costo adicional:
\[
\gamma\left[V_{p_0}(\mu^S)-V_{p_0}(\mu^F)\right]
\ge
\frac{\Delta c}{\Delta p},
\qquad
\Delta c=c_2-c_1,
\quad
\Delta p=p_2-p_1.
\]
Con los parámetros base,
\[
\frac{\Delta c}{\Delta p}
=\frac{0.65-0.05}{0.80-0.35}
=1.333\overline{3}.
\]

\section{Método numérico}

La ecuación de Bellman se resuelve por value iteration en una grilla uniforme
\(\mu\in[\varepsilon,1-\varepsilon]\), con \(\varepsilon=10^{-5}\) y 1000 puntos.
Las evaluaciones fuera de la grilla exacta, como \(V(\mu^S)\) y \(V(\mu^F)\), se
hacen por interpolación lineal. Para cada \(p_0\) y cada modo de demanda se
guardan:
\[
V_{p_0}(\mu), \quad
x^*_{p_0}(\mu), \quad
G_{p_0}(\mu)
=
\gamma[V_{p_0}(\mu^S)-V_{p_0}(\mu^F)],
\quad
D(\mu;p_0).
\]
También se simulan trayectorias desde \(\mu_{\mathrm{init}}=0.5\). En cada
periodo se registra la creencia, la demanda, si el usuario elige al vendedor, el
producto elegido, el resultado binario y la ganancia descontada.

\section{Resultados principales}

\subsection{\texorpdfstring{Regímenes de \(p_0\)}{Regimenes de p0}}

Los resultados se organizan en tres regímenes:
\[
\texttt{below\_p1}:p_0<p_1,
\qquad
\texttt{between\_p1\_p2}:p_1\le p_0\le p_2,
\qquad
\texttt{above\_p2}:p_0>p_2.
\]
En los numéricos, todos los casos no triviales convergen con residual Bellman de
orden \(10^{-9}\). Cuando \(p_0>p_2\), la demanda es cero, la función de valor es
cero y el vendedor no usa el producto 2.

\subsection{Demanda por umbral de media}

En \texttt{mean\_threshold}, la política muestra una banda interior para varios
valores intermedios de \(p_0\). A medida que \(p_0\) aumenta, el umbral de demanda
sube y la banda donde el producto 2 es óptimo se desplaza hacia creencias más
altas. La fracción de la grilla donde se usa producto 2 cae de \(0.947\) en
\(p_0=0.40\) a \(0.221\) en \(p_0=0.75\).

Un punto importante es que la política óptima puede recomendar producto 2 para
creencias altas, aunque la simulación desde \(\mu_{\mathrm{init}}=0.5\) no llegue
a esas regiones. Por ejemplo, con \(p_0=0.65\), la región de producto 2 empieza
aproximadamente en \(\mu=0.467\), pero la demanda del usuario empieza recién en
\(\mu_0(0.65)=0.667\). Como la trayectoria empieza en \(0.5\), no hay demanda,
no se aprende y la ganancia simulada es cero.

\subsection{Thompson sampling sobre tipo}

En \texttt{ts\_type}, dentro del régimen intermedio \(p_1<p_0\le p_2\), la
demanda esperada es \(D(\mu;p_0)=\mu\). Por lo tanto, la política y la función de
valor son esencialmente constantes en \(p_0\) dentro de ese régimen. En los
resultados, para \(p_0=0.40,\ldots,0.80\), la banda de producto 2 es
aproximadamente
\[
[0.005,\;0.999],
\]
con fracción de grilla \(0.994\) y \(V(\mu_{\mathrm{init}})\approx 32.704\).
Esta regla conserva demanda positiva aun cuando \(\bar p(\mu_{\mathrm{init}})\)
está por debajo de \(p_0\), porque la exploración viene del muestreo del tipo.

\begin{table}[H]
\centering
\small
\begin{tabular}{llrrrrr}
\toprule
\(p_0\) & Modo & \(V(\mu_{\mathrm{init}})\) & Frac. \(x=2\) &
Demanda sim. & Uso \(x=2\) sim. & Profit desc. sim. \\
\midrule
0.30 & mean & 47.500 & 0.000 & 1.000 & 0.000 & 47.498 \\
0.50 & mean & 26.440 & 0.820 & 0.735 & 0.402 & 21.164 \\
0.65 & mean & 0.000 & 0.532 & 0.000 & 0.000 & 0.000 \\
0.90 & mean & 0.000 & 0.000 & 0.000 & 0.000 & 0.000 \\
\midrule
0.30 & TS & 47.500 & 0.000 & 1.000 & 0.000 & 47.498 \\
0.50 & TS & 32.704 & 0.994 & 0.978 & 0.534 & 24.755 \\
0.65 & TS & 32.704 & 0.994 & 0.960 & 0.529 & 24.721 \\
0.90 & TS & 0.000 & 0.000 & 0.000 & 0.000 & 0.000 \\
\bottomrule
\end{tabular}
\caption{Resultados seleccionados desde \(\mu_{\mathrm{init}}=0.5\). ``mean''
corresponde a \texttt{mean\_threshold}; ``TS'' corresponde a
\texttt{ts\_type}.}
\label{tab:selected-results}
\end{table}

\section{Figuras}

La Figura~\ref{fig:comparative-statics} resume los estáticos comparativos. Las
líneas verticales punteadas marcan \(p_1=0.35\) y \(p_2=0.80\).

\begin{figure}[H]
\centering
\includegraphics[width=\textwidth]{figures_hidden_type/comparative_statics.png}
\caption{Estáticos comparativos en \(p_0\), para ambos modos de demanda.}
\label{fig:comparative-statics}
\end{figure}

La Figura~\ref{fig:mean-p050} muestra un caso interior en el modo de umbral de
media. La demanda salta en \(\mu_0(0.50)=1/3\), mientras que la banda óptima de
producto 2 empieza cerca de \(\mu=0.179\).

\begin{figure}[H]
\centering
\includegraphics[width=\textwidth]{figures_hidden_type/solution_mean_threshold_p0_0p50.png}
\caption{Solución para \texttt{mean\_threshold} con \(p_0=0.50\).}
\label{fig:mean-p050}
\end{figure}

La Figura~\ref{fig:mean-path-p065} ilustra el caso \(p_0=0.65\). Aunque existe
una región de creencias altas donde el producto 2 sería óptimo, la trayectoria
desde \(\mu_{\mathrm{init}}=0.5\) no genera demanda porque el umbral del usuario
es \(\mu_0=2/3\).

\begin{figure}[H]
\centering
\includegraphics[width=\textwidth]{figures_hidden_type/sample_path_mean_threshold_p0_0p65.png}
\caption{Trayectoria simulada para \texttt{mean\_threshold} con \(p_0=0.65\).}
\label{fig:mean-path-p065}
\end{figure}

La Figura~\ref{fig:ts-p065} contrasta el mismo \(p_0=0.65\) bajo
\texttt{ts\_type}. Como la demanda esperada es \(\mu\), la trayectoria mantiene
posibilidad de observaciones y aprendizaje desde \(\mu_{\mathrm{init}}=0.5\).

\begin{figure}[H]
\centering
\includegraphics[width=\textwidth]{figures_hidden_type/solution_ts_type_p0_0p65.png}
\caption{Solución para \texttt{ts\_type} con \(p_0=0.65\).}
\label{fig:ts-p065}
\end{figure}

\section{Conclusiones}

El modelo escalar produce una estructura más limpia que el estado Beta
\((S,F)\). En particular:
\begin{enumerate}
    \item El estado relevante es una sola creencia \(\mu\), lo que permite ver
    directamente cómo \(p_0\) desplaza la demanda y la política.
    \item Bajo \texttt{mean\_threshold}, \(p_0\) afecta fuertemente tanto la
    demanda como el valor: si \(\mu_{\mathrm{init}}\) queda por debajo del
    umbral de demanda, la economía se queda sin aprendizaje.
    \item Bajo \texttt{ts\_type}, \(p_0\) sólo determina el régimen. Dentro de
    \(p_1<p_0\le p_2\), la Bellman no cambia con \(p_0\), por lo que las
    políticas son casi planas en los estáticos comparativos.
    \item La variable clave para la elección del producto es el gap de
    continuación \(G_{p_0}(\mu)\). El producto 2 se usa cuando ese gap supera
    \(\Delta c/\Delta p\).
\end{enumerate}

\end{document}
